\documentclass[../otf-unofficial-guide]{subfiles}
\begin{document}

\section{数式への和文挿入}

{\upLaTeX*}は、数式中で明朝体を使用するための\cmd{mathmc}、ゴシック体を使用するための\cmd{mathgt}を提供しています。これに加えて、\pkg{otf}の\opt{deluxe}下では、丸ゴシック体を使用するための\cmd{mathmg}が提供されます。ただし、多ウェイト化はされていないため、\cmd{mathmc}および\cmd{mathgt}は{\upLaTeX*}デフォルトのウェイトが使用されます。

以下のように使用されます。

\begin{demosource}[0.5]
  {\textbackslash}[\\
    \tabto{2zw}{\textbackslash}mathmc\{\jpincmd{明朝}\}\\
    \tabto{2zw}{\textbackslash}mathgt\{\jpincmd{ゴシック}\}\\
    \tabto{2zw}{\textbackslash}mathmg\{\jpincmd{丸ゴシック}\}\\
  {\textbackslash}]
  \tcblower
  \[\mathmc{明朝}\mathgt{ゴシック}\mathmg{丸ゴシック}\]
\end{demosource}

\opt{disablejfam}%
\sidenote{\opt{disablejfam}は数式モード内で和文を利用しないようにするオプションです。}%
が有効になっているとき、\cmd{mathmc}、\cmd{mathgt}は無効なコマンドになりますが、\cmd{mathmg}は無効にならないことに注意してください。

\end{document}
